\documentclass[11pt]{article}

\usepackage[T1]{fontenc}
\usepackage[utf8]{inputenc}
\usepackage{geometry}
\usepackage{hyperref}
\usepackage{xcolor}
\usepackage{enumitem}
\usepackage{fancyvrb}
\usepackage{titlesec}

\geometry{margin=1in}
\hypersetup{
  colorlinks=true,
  linkcolor=blue!60!black,
  urlcolor=blue!60!black,
  pdftitle={ArmaMCP Initial Setup and Activation Guide},
  pdfauthor={ArmaMCP}
}

\setlist[itemize]{leftmargin=1.4em}
\setlist[enumerate]{leftmargin=1.6em}
\titleformat{\section}{\Large\bfseries}{\thesection}{0.7em}{}
\titleformat{\subsection}{\large\bfseries}{\thesubsection}{0.7em}{}

\newcommand{\code}[1]{\texttt{#1}}
\newcommand{\projectroot}{\code{/path/to/arma-mcp}}

\DefineVerbatimEnvironment{commandblock}{Verbatim}{
  fontsize=\small,
  frame=single,
  framerule=0.3pt,
  framesep=6pt
}

\title{ArmaMCP Initial Setup and Activation Guide}
\author{ArmaMCP MVP}
\date{June 7, 2026}

\begin{document}
\maketitle

\section{Purpose}

ArmaMCP connects Codex to Arma 3 Eden Editor through a local Model Context Protocol
(MCP) sidecar, a localhost HTTP bridge, a native Arma extension, and a HEMTT-built
SQF addon.

\begin{commandblock}
Codex MCP client
  -> TypeScript stdio sidecar
  -> 127.0.0.1 HTTP bridge
  -> ArmaMCP_x64 native extension
  -> SQF Eden addon
\end{commandblock}

The MVP activation path lets Codex ping the sidecar, request Eden selection
snapshots, generate a dry-run checkpoint plan, queue a reviewed plan, and inspect
bridge events from Eden.

\section{Safety Model}

ArmaMCP is local-first by design.

\begin{itemize}
  \item The HTTP bridge binds to \code{127.0.0.1} by default.
  \item Bridge endpoints under \code{/bridge/*} require \code{Authorization: Bearer <ARMA\_MCP\_TOKEN>}.
  \item The sidecar token is supplied with the \code{ARMA\_MCP\_TOKEN} environment variable.
  \item The Arma extension reads the matching token from \code{ArmaMCP.ini} next to the loaded extension.
  \item No OpenAI keys or durable secrets should be placed in SQF, PBOs, mission files,
        extension source, compiled binaries, or committed config files.
  \item The MVP exposes no raw SQF execution tool, no shell execution, no remote execution,
        and no public multiplayer server control.
  \item Queued editor plans are restricted to structured \code{createObject} and
        \code{createMarker} operations.
\end{itemize}

\section{Prerequisites}

\begin{itemize}
  \item Node.js and npm for the TypeScript sidecar.
  \item HEMTT on \code{PATH} for the Arma addon build.
  \item CMake and a C++17 compiler for the native Linux extension build.
  \item \code{x86\_64-w64-mingw32-g++} if building the Windows/Proton DLL locally.
  \item Arma 3 installed. The local Arma 3 Tools path assumed by the project is:
\end{itemize}

\begin{commandblock}
/mnt/game_one/SteamLibrary/steamapps/common/Arma 3 Tools/
\end{commandblock}

\section{Repository Layout}

The relevant project paths are:

\begin{itemize}
  \item \projectroot{}: repository root.
  \item \code{sidecar/}: TypeScript MCP server and HTTP bridge.
  \item \code{addons/main/}: HEMTT addon and Eden SQF functions.
  \item \code{extension/}: native \code{ArmaMCP\_x64} extension source.
  \item \code{scripts/build-extension.sh}: native extension build helper.
  \item \code{scripts/validate.sh}: full local validation helper.
\end{itemize}

\section{Build The Sidecar}

From the repository root:

\begin{commandblock}
cd /path/to/arma-mcp/sidecar
npm install
npm run build
\end{commandblock}

To run the sidecar tests:

\begin{commandblock}
cd /path/to/arma-mcp/sidecar
npm run validate
\end{commandblock}

The test suite opens temporary \code{127.0.0.1} ports. If a sandbox blocks local
binding, the HTTP tests can fail with \code{listen EPERM}; run them in an
environment that permits loopback binding.

\section{Choose A Local Bridge Token}

Pick a local token and keep it out of git. For example:

\begin{commandblock}
export ARMA_MCP_TOKEN='replace-with-a-local-random-token'
\end{commandblock}

Use the same token for:

\begin{itemize}
  \item the sidecar process;
  \item the Codex MCP server configuration;
  \item the local mod folder's \code{ArmaMCP.ini}.
\end{itemize}

For quick local-only smoke tests, documentation examples use \code{dev-token}, but
that value is only a placeholder.

\section{Configure Codex MCP}

First build the sidecar:

\begin{commandblock}
cd /path/to/arma-mcp/sidecar
npm run build
\end{commandblock}

Add this server entry to Codex \code{config.toml}, replacing the token with your
local value:

\begin{commandblock}
[mcp_servers.arma-mcp]
command = "node"
args = ["/path/to/arma-mcp/sidecar/dist/index.js"]
env = { ARMA_MCP_TOKEN = "replace-with-a-local-random-token" }
startup_timeout_sec = 10
tool_timeout_sec = 60
\end{commandblock}

If Arma is already polling a separately started bridge, configure Codex to reuse
that bridge instead of binding the HTTP listener again:

\begin{commandblock}
[mcp_servers.arma-mcp]
command = "node"
args = ["/path/to/arma-mcp/sidecar/dist/index.js"]
env = {
  ARMA_MCP_TOKEN = "replace-with-a-local-random-token",
  ARMA_MCP_SKIP_HTTP_LISTEN = "1"
}
startup_timeout_sec = 10
tool_timeout_sec = 60
\end{commandblock}

Alternatively, register it with the Codex CLI:

\begin{commandblock}
codex mcp add arma-mcp \
  --env ARMA_MCP_TOKEN=replace-with-a-local-random-token \
  -- node /path/to/arma-mcp/sidecar/dist/index.js
\end{commandblock}

Reuse-mode helper form:

\begin{commandblock}
codex mcp add arma-mcp \
  --env ARMA_MCP_TOKEN=replace-with-a-local-random-token \
  --env ARMA_MCP_SKIP_HTTP_LISTEN=1 \
  -- node /path/to/arma-mcp/sidecar/dist/index.js
\end{commandblock}

After starting Codex, inspect MCP servers with:

\begin{commandblock}
/mcp
\end{commandblock}

Expected result: \code{arma-mcp} is listed and its tools are available.

\section{Run The Bridge Without Codex}

For bridge-only development:

\begin{commandblock}
cd /path/to/arma-mcp/sidecar
ARMA_MCP_TOKEN=dev-token npm run dev:http
\end{commandblock}

This starts only the HTTP bridge. Normal MCP mode should be launched by Codex and
keeps MCP traffic on stdout while logs go to stderr.

For local development of a Codex stdio process that reuses the running bridge:

\begin{commandblock}
ARMA_MCP_TOKEN=dev-token npm run dev:stdio-existing
\end{commandblock}

\section{Build The Native Extension}

From the repository root:

\begin{commandblock}
cd /path/to/arma-mcp
./scripts/build-extension.sh
\end{commandblock}

Expected generated files:

\begin{itemize}
  \item \code{extension/build/ArmaMCP\_x64.so}
  \item \code{extension/build/ArmaMCP\_x64.dll}, when MinGW is available
\end{itemize}

The generated build directory is ignored by git.

\subsection{Proton Path}

For Arma 3 running under Proton:

\begin{commandblock}
Arma 3 Windows executable under Proton
  -> loads ArmaMCP_x64.dll
  -> talks to native Linux sidecar over 127.0.0.1
\end{commandblock}

Use the Windows DLL for client-side Eden testing under Proton. The native Linux
\code{.so} is useful for local smoke testing and future native server work, but it
is not the first Proton client path.

\section{Build The HEMTT Addon}

From the repository root:

\begin{commandblock}
cd /path/to/arma-mcp
hemtt build
\end{commandblock}

For a signed local release:

\begin{commandblock}
./scripts/release-signed.sh
\end{commandblock}

This runs \code{hemtt release}, checks that the PBO has a matching \code{.bisign},
checks that HEMTT emitted the public \code{.bikey}, and verifies the signature
with \code{hemtt utils verify}.

For a development deploy:

\begin{commandblock}
hemtt dev
\end{commandblock}

\code{hemtt dev} builds the addon and then tries to deploy to the local Arma 3
installation. If the Arma installation path is not writable, the build can succeed
while deployment fails.

\section{Configure The Arma Extension}

Create \code{ArmaMCP.ini} next to \code{ArmaMCP\_x64.dll} in the local mod folder:

\begin{commandblock}
@ArmaMCP/
  ArmaMCP_x64.dll
  ArmaMCP.ini
  addons/amcp_main.pbo
\end{commandblock}

The config file format is:

For a simple local smoke test:

\begin{commandblock}
host=127.0.0.1
port=38473
token=dev-token
\end{commandblock}

The important rule is that the extension config file and sidecar use the same token.
Environment variables \code{ARMA\_MCP\_HOST}, \code{ARMA\_MCP\_PORT}, and
\code{ARMA\_MCP\_TOKEN} remain available as fallback settings when
\code{ArmaMCP.ini} is absent.

\section{Initial Eden Activation Test}

\begin{enumerate}
  \item Build the sidecar with \code{npm run build}.
  \item Configure Codex MCP with the \code{arma-mcp} sidecar.
  \item Build the native extension with \code{./scripts/build-extension.sh}.
  \item Build the addon with \code{hemtt build} or deploy with \code{hemtt dev}.
  \item Make \code{ArmaMCP\_x64.dll} available where Arma can load it.
  \item Create \code{ArmaMCP.ini} beside the DLL with the same token as the sidecar.
  \item Launch Arma with the addon loaded and open Eden.
  \item Place a simple object or game logic and select it.
  \item In Codex, call \code{arma\_ping}.
  \item Call \code{arma\_request\_editor\_snapshot}.
  \item Wait for Eden to poll, then call \code{arma\_get\_editor\_snapshot}.
  \item Call \code{arma\_generate\_checkpoint\_plan}.
  \item Review the returned dry-run JSON.
  \item Set \code{dryRun=false} only after review, then call \code{arma\_queue\_apply\_plan}.
  \item Call \code{arma\_get\_bridge\_events} and inspect the result.
  \item Confirm Eden creates objects around the selected anchor.
  \item Press undo in Eden to confirm the operation rolls back if
        \code{collect3DENHistory} grouped changes successfully.
\end{enumerate}

\section{Available MCP Tools}

The current sidecar exposes dotted MCP tools for bridge diagnostics, Eden reads
and writes, catalog work, camera/screenshot capture, data-only composition
helpers, terrain sampling, and procedural generators. Useful starting tools are:

\begin{itemize}
  \item \code{arma.ping}: check MCP stdio process health without waiting on Eden.
  \item \code{arma.bridge.diagnostics}: report stdio mode, listener ownership,
        bridge reachability, recent bridge activity, and recent bridge errors.
  \item \code{arma.bridge.get\_status}: check pending actions/results and the
        latest Eden bridge contact.
  \item \code{arma.bridge.get\_capabilities}: inspect addon/runtime capabilities.
  \item \code{arma.eden.get\_status}, \code{arma.eden.get\_selection}, and
        \code{arma.eden.list\_entities}: inspect live Eden state before planning.
  \item \code{arma.eden.batch}, \code{arma.eden.validate\_plan}, and
        \code{arma.eden.apply\_composition}: dry-run and apply reviewed changes.
  \item \code{arma.catalog.scanStart}, \code{arma.catalog.scanStatus}, and
        \code{arma.catalog.scanPoll}: start, inspect, and advance bounded
        catalog scan work.
  \item \code{arma.catalog.scanCancel}: cancel a running catalog scan.
  \item \code{arma.catalog.scanFinalize}: finish a completed bounded scan.
  \item \code{arma.catalog.scanRepair}: repair stale scan bookkeeping.
  \item \code{arma\_visual\_inspect\_class}: discovery-friendly alias for
        \code{arma.visual.inspectClass}.
  \item \code{arma\_composition\_plan}: discovery-friendly alias for
        \code{arma.composition.plan}.
  \item \code{arma\_discovery} and \code{arma\_call}: list and call an
        allowlisted fallback surface when a Codex-managed session omits
        individual tools that direct stdio exposes.
\end{itemize}

Legacy MVP aliases remain available for compatibility:

\begin{itemize}
  \item \code{arma\_ping}
  \item \code{arma\_request\_editor\_snapshot}
  \item \code{arma\_get\_editor\_snapshot}
  \item \code{arma\_generate\_checkpoint\_plan}
  \item \code{arma\_queue\_apply\_plan}
  \item \code{arma\_get\_bridge\_events}
\end{itemize}

\section{Troubleshooting}

\subsection{Codex Does Not List \code{arma-mcp}}

\begin{itemize}
  \item Confirm \code{sidecar/dist/index.js} exists.
  \item Confirm the Codex MCP config path is absolute and points to this repo.
  \item Confirm \code{ARMA\_MCP\_TOKEN} is set in the MCP server environment.
  \item Check sidecar stderr logs.
  \item If direct stdio exposes tools that Codex-managed discovery omits,
        restart/reload the Codex MCP host; use \code{arma\_discovery} and
        \code{arma\_call} as the bounded fallback while the registry is stale.
  \item For high-level visual inspection, use \code{arma\_visual\_inspect\_class}
        when managed discovery omits the dotted tool name.
  \item For high-level composition planning, use \code{arma\_composition\_plan}
        when managed discovery omits the dotted tool name.
\end{itemize}

\subsection{\code{arma\_ping} Works But Eden Is Not Connected}

\begin{itemize}
  \item Run \code{arma.bridge.diagnostics} to confirm whether Codex stdio is
        reusing an existing bridge or owns the HTTP listener.
  \item Confirm Arma is running with the addon loaded.
  \item Confirm the native extension is available to Arma as \code{ArmaMCP\_x64.dll}
        for Proton/Windows.
  \item Confirm \code{ArmaMCP.ini} sits beside \code{ArmaMCP\_x64.dll}.
  \item Confirm \code{ArmaMCP.ini} uses the same token as the sidecar.
  \item Confirm the sidecar is listening on \code{127.0.0.1:38473}.
\end{itemize}

\subsection{HEMTT Build Or Deploy Fails}

\begin{itemize}
  \item Use \code{hemtt build} to validate config and SQF compilation.
  \item Use \code{hemtt dev} only when the Arma installation path is writable.
  \item Deployment failures can happen after a successful PBO build.
\end{itemize}

\subsection{Extension Build Fails}

\begin{itemize}
  \item Confirm \code{cmake} and a C++17 compiler are installed.
  \item Install \code{mingw-w64} if \code{ArmaMCP\_x64.dll} is needed locally.
  \item The Linux \code{.so} can build even if the Windows cross-compiler is missing.
\end{itemize}

\subsection{Plan Queue Is Rejected}

\begin{itemize}
  \item \code{arma\_queue\_apply\_plan} rejects \code{dryRun=true}.
  \item Unsupported operation types are rejected.
  \item The MVP accepts at most 50 operations.
  \item Only \code{createObject} and \code{createMarker} are accepted.
\end{itemize}

\section{Known MVP Limits}

\begin{itemize}
  \item Sidecar state is in-memory only.
  \item The addon assumes current Arma 3 support for \code{fromJSON} and \code{toJSON}.
  \item The extension reads token, host, and port from \code{ArmaMCP.ini}, with environment variables as fallback settings.
  \item Live class lookup, local catalog scans, terrain sampling, road probes, and
        broad spatial validation are implemented as review aids. Mesh parsing,
        exact collision validation, and fully faction-aware catalog selection are
        still future work.
  \item Marker creation still needs live Eden verification.
  \item BattleEye/public-server deployment, durable audit logs, and admin
        authority controls are future work.
\end{itemize}

\section{One-Command Validation}

After local prerequisites are installed:

\begin{commandblock}
cd /path/to/arma-mcp
./scripts/validate.sh
\end{commandblock}

This runs sidecar validation, extension build, HEMTT build, signed HEMTT release
verification, and git proof commands.

\end{document}
